The stochastic inverse problem augments the same logistic drift with
multiplicative process noise:
$$
  dX_t
  =
  rX_t\left(1-\frac{X_t}{K}\right)\,dt
  +
  \sigma X_t\, \circ dW_t .
$$
The stochastic PINN retains the same neural trajectory, data loss, initial
loss, and physics residual used by the deterministic PINN. The additional
identification of $\sigma$ comes from an Euler-Maruyama transition likelihood
on adjacent observed pairs. For $\Delta t_i=t_{i+1}-t_i$,
$$
  Y_{i+1}^{[j]}\mid Y_i^{[j]}
  \approx
  \mathcal N(m_i^{[j]},v_i^{[j]}),
$$
where
$$
  m_i^{[j]}
  =
  y_i^{[j]}
  +
  r y_i^{[j]}\left(1-\frac{y_i^{[j]}}{K}\right)\Delta t_i,
  \qquad
  v_i^{[j]}
  =
  \left(\sigma y_i^{[j]}\right)^2\Delta t_i+10^{-6}.
$$
Let $\Omega_{\mathrm{tr}}$ be the set of adjacent observed transitions. The
transition negative log-likelihood is
$$
  \mathcal L_{\mathrm{tr}}(r,K,\sigma)
  =
  \frac{1}{|\Omega_{\mathrm{tr}}|}
  \sum_{(i,j)\in\Omega_{\mathrm{tr}}}
  \frac{1}{2}
  \left[
    \log(2\pi v_i^{[j]})
    +
    \frac{(y_{i+1}^{[j]}-m_i^{[j]})^2}{v_i^{[j]}}
  \right].
$$
Operationally, this loss is computed one observed step at a time. For each
complete adjacent pair, the current value $y_i^{[j]}$ and time increment
$\Delta t_i$ determine the Euler-Maruyama mean $m_i^{[j]}$ and variance
$v_i^{[j]}$. The difference $y_{i+1}^{[j]}-m_i^{[j]}$ is the one-step
innovation around the logistic drift. The Gaussian negative log-likelihood
then combines two penalties: a scaled squared innovation
$(y_{i+1}^{[j]}-m_i^{[j]})^2/v_i^{[j]}$, which penalizes transitions that are
far from the predicted mean, and a $\log(2\pi v_i^{[j]})$ term, which penalizes
unnecessarily large transition variance. The reported transition NLL is the
average of these per-transition contributions over $\Omega_{\mathrm{tr}}$.
Thus a smaller value means that the observed kefir growth increments are more
plausible under the stochastic logistic dynamics. The parameter $\sigma$ is
identified through this term because it controls $v_i^{[j]}$: if $\sigma$ is
too small, ordinary short-time fluctuations are assigned very low probability,
whereas if $\sigma$ is too large, the log-variance part of the likelihood
penalizes the excess uncertainty.

The total stochastic objective is
$$
  \mathcal L_{\mathrm{stoPINN}}
  =
  \mathcal L_{\mathrm{data}}
  +
  \lambda_0\mathcal L_{\mathrm{init}}
  +
  \lambda_{\mathrm{phys}}\mathcal L_{\mathrm{phys}}
  +
  \lambda_{\mathrm{tr}}\mathcal L_{\mathrm{tr}},
$$
with $\lambda_0=1$, $\lambda_{\mathrm{phys}}=1$, and
$\lambda_{\mathrm{tr}}=0.1$. The transition likelihood is evaluated on the
original response scale and the original time increments, while the
collocation residual is evaluated in normalized time. Training used Adam for
3000 epochs with learning rate $5\times10^{-3}$.

The stochastic PINN/SDE estimation procedure extends the deterministic loop as
follows:
\begin{enumerate}
  \item Construct the same observation losses and collocation physics loss used
        by the deterministic PINN.
  \item Build a transition data set from adjacent observations:
        $(y_i^{[j]},y_{i+1}^{[j]},\Delta t_i)$ is included only when both
        endpoints are observed.
  \item Optimize three raw physical variables $a,b,c$, transformed to
        $r>0$, $K>K_{\min}$, and $\sigma>0$ with softplus maps.
  \item For every adjacent transition, compute the Euler-Maruyama conditional
        mean $m_i^{[j]}$ and variance $v_i^{[j]}$ implied by the current
        $r$, $K$, and $\sigma$.
  \item Compute the Gaussian transition negative log-likelihood
        $\mathcal L_{\mathrm{tr}}$. This term estimates the diffusion
        intensity because $\sigma$ controls the transition variance.
  \item Add the weighted transition likelihood to the PINN data, initial, and
        physics losses, then backpropagate through the full objective to update
        $(\theta,a,b,c)$.
  \item After training, report
        $\widehat r=r(\widehat a)$,
        $\widehat K=K(\widehat b)$, and
        $\widehat\sigma=\sigma(\widehat c)$.
\end{enumerate}

The stochastic fit uses two complementary sources of information. The PINN
terms estimate a smooth population drift trajectory and keep that trajectory
close to logistic growth. The transition likelihood estimates how much
short-time variability remains around the logistic drift. Consequently,
$\sigma$ is identified from the observed increments, not from the pointwise
distance between the PINN curve and the data.

After training, the stochastic logistic SDE was simulated with
Euler-Maruyama. Each trial-specific simulation starts from its observed initial
value $y_0^{[j]}$. The reported stochastic mean and central 90 percent
predictive interval use 5000 Monte Carlo paths. In the numerical table,
``Stochastic Logistic PINN Drift'' refers to the fitted neural trajectory
$u_\theta$, whereas ``Stochastic Logistic SDE Mean'' refers to the Monte Carlo
mean of the fitted SDE.